The $S_5$ classification suggests the natural converse to
Theorem~\ref{thm:sp-extension-tiler}: perhaps every extension tiler is
series-parallel.  This converse is false in the next rank.

\begin{theorem}[A non-series-parallel extension tiler in $S_6$]
\label{thm:s6-counterexample}
Let $P$ be the six-element poset with cover relations
\[
        0<1,\qquad 1<2,\qquad 1<3,\qquad
        4<3,\qquad 4<5,\qquad 5<2 .
\]
Then $P$ is not series-parallel, but $L(P)$ tiles $S_6$ by left
translates.
\end{theorem}

\begin{proof}
The poset contains induced copies of the forbidden $N$-poset; for example,
the induced subposet on $\{1,2,3,5\}$ has relations
\[
        1<2,\qquad 1<3,\qquad 5<2,
\]
which is isomorphic to $N$ under
\[
        a\mapsto 5,\qquad b\mapsto 2,\qquad
        c\mapsto 1,\qquad d\mapsto 3.
\]
Hence $P$ is not series-parallel.

We now give a compact multiplier set.  Let permutations act on the labels
of words, so that
\[
        g(w_1,\ldots,w_6)=(g(w_1),\ldots,g(w_6)).
\]
Define
\[
H=\langle (2\,3)(4\,5),\ (0\,1)(3\,5)\rangle,
\]
\[
K=\langle (2\,3)(4\,5),\ (0\,4\,5)(1\,3\,2)\rangle,
\]
and let $\tau=(3\,4)$.  Put
\[
        A=HK\sqcup H\tau K.
\]
A direct computation gives
\[
        |L(P)|=15,\qquad |H|=8,\qquad |K|=6,
\]
\[
        |HK|=|H\tau K|=24,\qquad HK\cap H\tau K=\emptyset.
\]
Thus $|A|=48$, which is the required number of tiles since
$48\cdot 15=720=|S_6|$.

It remains only to check disjointness.  The verification is finite:
compute the $15$ linear extensions of $P$, form all words $a\sigma$ with
$a\in A$ and $\sigma\in L(P)$, and count multiplicities.  The resulting
$720$ words are all distinct.  Hence
\[
        S_6=\bigsqcup_{a\in A} aL(P),
\]
so $P$ is an extension tiler.
\end{proof}

The verification in the last paragraph is intentionally small: it uses
only the two generated subgroups $H,K$ and the element $\tau$, rather than
a list of $48$ multipliers.  The supplementary script
\[
\texttt{scripts/verify\_s6\_p038\_biset\_counterexample.py}
\]
reproduces the calculation and also confirms that this poset contains
four induced copies of $N$.

\begin{corollary}
\label{cor:extension-tiler-converse-false}
The converse of Theorem~\ref{thm:sp-extension-tiler} is false.
\end{corollary}

\begin{corollary}[Coxeter-complex realization]
\label{cor:s6-counterexample-coxeter-realization}
Let $P$ be the six-element poset of
Theorem~\ref{thm:s6-counterexample}.  In the type $A_5$ Coxeter complex,
whose chambers are indexed by the elements of $S_6$, the chamber set
$L(P)$ is a $15$-chamber Coxeter-pure piece.  The $48$ left translates
\[
        \{aL(P): a\in A\}
\]
from Theorem~\ref{thm:s6-counterexample} form a disjoint partition of the
$720$ Coxeter chambers.  Equivalently, the verified $P038$ factorization
realizes a type $A_5$ Coxeter-pure group-translate tiling by $48$
translates of the $15$-chamber poset cone $C(P)=L(P)$.
\end{corollary}

\begin{proof}
Index the chambers of the type $A_5$ Coxeter complex by words
$w\in S_6$.  The chamber graph is the adjacent-transposition graph on
words.  Relabeling every entry of a word by an element of $S_6$ sends
adjacent words to adjacent words, and therefore preserves the Coxeter
chamber structure.

Theorem~\ref{thm:s6-counterexample} gives the exact factorization
\[
        S_6=\bigsqcup_{a\in A} aL(P),
\]
where $|L(P)|=15$ and $|A|=48$.  Replacing each word by the corresponding
Coxeter chamber gives a disjoint cover of all $720$ type $A_5$ chambers by
$48$ left translates of the $15$-chamber set $L(P)$.  Since each translate
is obtained by a Coxeter-group action, the pieces are congruent as
Coxeter-chamber unions.
\end{proof}

\begin{remark}
The multiplier set above has visible two-sided structure: it is a union of
two $(H,K)$-double cosets, with $H\cong D_4$ and $K\cong S_3$.  This
suggests that non-series-parallel extension tilers in $S_6$ are not merely
sporadic exact-cover artefacts; at least some arise from structured
group-theoretic complements.
\end{remark}
